\section{Ablation Study: Information Scarcity and Knowledge Contribution}
\label{sec:ablation}

To rigorously evaluate the contribution of Dynamic Knowledge Injection (DKI) and test the hypothesis of \textit{information saturation}, we conducted a systematic ablation study on the In-Hospital Mortality (IHM) task. We defined four clinical scenarios with decreasing levels of data density:
\begin{itemize}
    \item \textbf{Full (55 features)}: The complete aligned benchmark including all vitals, hemodynamics, laboratory values, and clinical interventions.
    \item \textbf{No Treatments (51 features)}: Removal of treatment-related variables (vasopressors, fluids, ventilation), simulating a purely diagnostic setting.
    \item \textbf{Clinical Core (15 features)}: Restricted to vital signs and essential SOFA-related laboratory values.
    \item \textbf{Emergency (10 features)}: A minimal set simulating early admission, including only vitals and Arterial Blood Gas (ABG) analysis.
\end{itemize}

In this phase, we introduced \textit{Deep Gated Injection} (DGI), where medical knowledge is injected into every layer of the Transformer via a contextual gating mechanism, ensuring that semantic signals persist through the deep feature extraction process.

\begin{table}[ht]
\centering
\caption{Ablation results for IHM. DGI (Deep Gated Injection) models utilize multi-layer semantic cross-attention. Best results per scenario are highlighted in bold.}
\label{tab:ablation_results}
\begin{tabular}{llccr}
\toprule
\textbf{Scenario} & \textbf{Model} & \textbf{AUROC} & \textbf{AUPRC} & \textbf{Note} \\
\midrule
\multirow{3}{*}{Full (55)} & Vanilla & \textbf{0.9063} & \textbf{0.5956} & Baseline \\
 & Grounded DKI (Input) & 0.9032 & 0.5865 & Single-layer \\
 & Deep DKI (DGI) & 0.9027 & 0.5851 & Multi-layer \\
\midrule
\multirow{3}{*}{\shortstack[l]{No Treatments\\(51)}} & Vanilla & 0.8972 & 0.5869 & \\
 & Grounded DKI (Input) & \textbf{0.9025} & 0.5880 & \\
 & Deep DKI (DGI) & 0.9017 & \textbf{0.5976} & \textit{State-of-the-Art} \\
\midrule
\multirow{2}{*}{Core (15)} & Vanilla & \textbf{0.8892} & \textbf{0.5632} & \\
 & DKI Adaptive & 0.8838 & 0.5627 & Input-only \\
\midrule
\multirow{2}{*}{Emergency (10)} & Vanilla & 0.8313 & 0.4553 & \\
 & DKI Adaptive & \textbf{0.8315} & \textbf{0.4619} & Input-only \\
\bottomrule
\end{tabular}
\end{table}

\subsection{Analysis of Findings}

\paragraph{Deep Semantic Persistence.}
The most significant finding is observed in the \textit{No Treatments} scenario, where the Deep Gated Injection (DGI) model achieved the highest AUPRC recorded in our experiments ($0.5976$). This represents a notable improvement over the Vanilla model ($0.5869$) and the input-level DKI ($0.5880$). By integrating knowledge at every abstraction level, DGI prevents the semantic signal from being overshadowed by numerical patterns, leading to more precise positive class identification in the absence of treatment data.

\paragraph{Knowledge as a Compensatory Mechanism.}
The consistent performance boost in AUPRC within information-sparse scenarios (\textit{No Treatments} and \textit{Emergency}) confirms that DKI acts as an effective \textit{relational prior}. It allows the model to reconstruct complex clinical dependencies that are otherwise unobservable when only a subset of features is available.

\paragraph{The Saturation Threshold.}
In the \textit{Full} dataset scenario, the superiority of the Vanilla model suggests an \textit{information saturation} effect. With 55 features, the Transformer's self-attention mechanism is already capable of learning near-optimal representations. In this high-density regime, additional knowledge injection may introduce minor architectural overhead or redundant signals that do not translate into further performance gains.
